\documentclass[aspectratio=169,professionalfonts]{beamer}
\usetheme{Madrid}
\usecolortheme{default}
\usepackage{amsmath,amssymb}
\usepackage[T1]{fontenc}
\title{STM-GARSOS Processing: Detailed Methods}
\author{STM-GARSOS workflow}
\date{\today}
\begin{document}
\begin{frame}\titlepage\end{frame}
\begin{frame}{Reconstruction objective}
\begin{itemize}
\item Dynamic image model: $X_t(x,y)=\sum_{\ell=1}^{L}{\rm ST}_{t,\ell}(x,y)C_\ell(x,y)$.
\item Unknown: C\_l(x,y), an [Nx x Ny x L] coefficient volume.
\item Measured data: multicoil radial samples for each temporal frame.
\item Goal: fit the measured radial data while restricting temporal variation to the STM subspace.
\end{itemize}
\end{frame}
\begin{frame}{Entry point and inputs}
\begin{itemize}
\item main\_STM\_GARSOS loads input k-space identifiers and scan metadata.
\item A vendor-reconstructed reference volume defines target geometry.
\item prepareGARSOSSTMData creates shared prepared data and calibration metadata.
\item Historical defaults: 15 spokes/frame, 32 x 32 calibration square, L = 7.
\end{itemize}
\end{frame}
\begin{frame}{Prepared trajectory and data}
\begin{itemize}
\item sampledData: [nReadout x nPartitions x nSpokes, nCoils].
\item kSpaceLocations: flattened [kx, ky, kz] rows.
\item spokeIndex: one identifier per flattened sample.
\item DCF: one radial density-compensation value per sample.
\item Coordinates use normalized cycles/FOV-like units.
\end{itemize}
\end{frame}
\begin{frame}{Trajectory construction}
\begin{itemize}
\item $k_x=\rho\cos(\phi)$ and $k_y=\rho\sin(\phi)$.
\item rho is the centered readout coordinate normalized by readout length.
\item phi follows the GAR-SOS spoke-angle convention.
\item kz is generated from the partition index and acquisition geometry.
\item The arrays are flattened in readout/partition/spoke order.
\end{itemize}
\end{frame}
\begin{frame}{Readout processing and whitening}
\begin{itemize}
\item Readout delay correction is estimated from selected calibration lines.
\item Readout Fourier processing is applied before the final flattening.
\item Partition padding is performed during preparation.
\item Noise covariance is estimated from noise acquisitions.
\item The whitening transform is applied across receive channels.
\end{itemize}
\end{frame}
\begin{frame}{Density compensation}
\begin{itemize}
\item The existing DCF is based on radial shell-sector area.
\item It is generated from the in-plane radius and readout increment.
\item Frame-specific DCF values are selected with the frame samples.
\item Calibration gridding applies DCF and additional reconstruction weights.
\item The later STM radial operator uses its prepared radial samples and NUFFT convention.
\end{itemize}
\end{frame}
\begin{frame}{Frame grouping}
\begin{itemize}
\item Unique spoke identifiers define the acquisition order.
\item nFrames = floor(nSpokes / nSpokesPerFrame).
\item Each frame contains all readout and partition samples for its spokes.
\item Incomplete trailing spokes are excluded.
\item For the baseline, 15 spokes/frame produced approximately 133 frames.
\end{itemize}
\end{frame}
\begin{frame}{Calibration reconstruction}
\begin{itemize}
\item reconstructGARFrameForSTM selects frameData, frameTrajectory, and frameDCF.
\item Coordinate transform and translation phase correction are applied.
\item Gaussian/spatial weighting and DCF scaling are applied.
\item GpuReconstructur20181008 grids every coil separately.
\item The class applies inverse FFT, resize/deapodization, and coil-profile accumulation.
\end{itemize}
\end{frame}
\begin{frame}{3-D grid and target geometry}
\begin{itemize}
\item Target image geometry: imageSize = [224, 224, 96].
\item Oversampled grid: gridSize = ceil(1.375 x imageSize) = [308, 308, 132].
\item Samples are deposited with a 3-D Kaiser-Bessel kernel.
\item ifftn is applied on the oversampled Cartesian grid.
\item shiftAndResize maps the result to the centered 224 x 224 x 96 grid.
\end{itemize}
\end{frame}
\begin{frame}{Calibration k-space extraction}
\begin{itemize}
\item Select reconstructedVolume(:, :, zSlice).
\item Apply fftshift(fft2(ifftshift(senseImage))).
\item Extract the centered calibration square with localCentralRegion.
\item Baseline calibrationSize = [32, 32].
\item Single-slice kCal has dimensions [32 x 32 x nFrames].
\end{itemize}
\end{frame}
\begin{frame}{STM computation}
\begin{itemize}
\item $\texttt{STM\_computation}$ receives $k_{\rm Cal}$, desired map dimensions, $L$, and options.
\item The calibration data encode temporal correlations across frames.
\item The algorithm estimates ST\_maps [Nx x Ny x Nt x L].
\item eigenValues provide a diagnostic of model rank/content.
\item The baseline used L = 7 because more modes appeared informative than in the Cartesian demo.
\end{itemize}
\end{frame}
\begin{frame}{Historical z-slice preparation}
\begin{itemize}
\item prepareGARSOSSliceDataForMCNUFFT restores [readout, partition, spoke, coil].
\item It applies the selected translation phase and coordinate transform.
\item Historical baseline: centered direct IFFT along the 58 partition entries.
\item The requested zIndex is selected from that 58-point result.
\item This differs from the later 58-to-96 gridded z-only correction.
\end{itemize}
\end{frame}
\begin{frame}{Radial STM data dimensions}
\begin{itemize}
\item After frame grouping: radialData [nReadout x nSpokesPerFrame x nCoil x nFrame].
\item kxy [nReadout x nSpokesPerFrame x nFrame].
\item The in-plane radial trajectory is retained.
\item The coil dimension is not combined before the STM solve.
\item Trajectory/data correspondence is validated before operator construction.
\end{itemize}
\end{frame}
\begin{frame}{Forward sensitivity-map model}
\begin{itemize}
\item Forward coil model: $d_c={\cal N}_{\rm radial}(S_c\odot X)$.
\item S\_c are forward receive sensitivities in the whitened channel basis.
\item Adjoint coil combination is $\sum_c S_c^* d_c$ after NUFFT adjoints.
\item This is distinct from calibration-only coil-combination profiles.
\item Map and dynamic-data channel ordering must match.
\end{itemize}
\end{frame}
\begin{frame}{Radial NUFFT operator}
\begin{itemize}
\item buildGARSOSRadialEncodingOperator creates one NUFFT plan per frame.
\item Forward: multiply by coil maps, then NUFFT each coil.
\item Adjoint: NUFFT-adjoint each coil, then sum conjugated maps.
\item Optional single precision converts trajectories, maps, data, and outputs.
\item The operator exposes forward, adjoint, and vectorized forms.
\end{itemize}
\end{frame}
\begin{frame}{STM operator composition}
\begin{itemize}
\item buildGARSOSSTMOperators composes STM synthesis with radial encoding.
\item $\mathcal{A}(C)=\mathcal{F}(\mathcal{S}(C))$.
\item $\mathcal{A}^H(y)=\mathcal{S}^H(\mathcal{F}^H(y))$.
\item The normal operator is applied without explicitly forming a matrix.
\item The coefficient vector has Nx x Ny x L entries.
\end{itemize}
\end{frame}
\begin{frame}{PCG solve}
\begin{itemize}
\item Right-hand side: $b=\mathcal{A}^H y$.
\item Unregularized system: $(\mathcal{A}^H\mathcal{A})C=b$.
\item The baseline used tolerance 1e-5 and up to 50 iterations.
\item Residual history, status flag, iteration count, and elapsed time are stored.
\item Preconditioning and Toeplitz support were preliminary rather than baseline assumptions.
\end{itemize}
\end{frame}
\begin{frame}{Optional Tikhonov solve}
\begin{itemize}
\item Regularized system: $(\mathcal{A}^H\mathcal{A}+\lambda I)C=b$.
\item The implementation can run a second PCG solve with lambda.
\item The effective lambda depends on data, map, NUFFT, and DCF scaling.
\item A numerical lambda should not be compared across differently scaled operators without calibration.
\end{itemize}
\end{frame}
\begin{frame}{Synthesis and outputs}
\begin{itemize}
\item $\widehat{X}_t=\sum_{\ell=1}^{L}{\rm ST}_{t,\ell}\widehat{C}_\ell$.
\item sense\_recon\_stm has dimensions [Nx x Ny x nFrame].
\item The result also contains coefficient images and PCG diagnostics.
\item The initial reported output contained approximately 133 reconstructed frames.
\end{itemize}
\end{frame}
\begin{frame}{Baseline versus later improvements}
\begin{itemize}
\item Baseline: direct 58-point partition IFFT for STM slice data.
\item Later work: forward sensitivity-map estimator and map-semantics audit.
\item Later work: coil-preserving gridded z-only path for the 96-slice geometry.
\item Later work: scaling, conditioning, coverage, regularization, and preconditioning analysis.
\item These improvements should be presented separately from the historical baseline.
\end{itemize}
\end{frame}
\begin{frame}{End-to-end summary}
\begin{itemize}
\item Raw GAR-SOS data -> corrected, whitened radial samples.
\item Trajectory + DCF + spoke IDs -> temporal frames.
\item 3-D gridding -> calibration Cartesian k-space -> central square.
\item STM calibration -> ST maps and eigenvalues.
\item Coil-resolved radial data + sensitivity maps -> STM/radial operator.
\item PCG -> spatial coefficients -> dynamic image series.
\end{itemize}
\end{frame}
\end{document}
